% Generated by scripts/assemble.py; corrections belong in source/corrections.json.
\examhead{80}{东北师范大学 2026 年数学分析试卷(回忆版)}{78}
% p079-b15; source PDF page 79
\blockanchor{p079-b15}
(注. 题目来自考生回忆, 部分表述或数据可能有误)\par
% p079-b16; source PDF page 79
\par\Needspace{5\baselineskip}
\blockanchor{p079-b16}
\majorquestion{一}计算题. 每题6分, 共30分.\par
% p079-b17; source PDF page 79
\blockanchor{p079-b17}
\question{1}求极限 $\lim_{n\to \infty}(2026\cos^2 n + \sin^2 n)^{\frac{1}{n}}$\par
% p079-b18; source PDF page 79
\blockanchor{p079-b18}
\question{2}求极限  $\lim_{x\to0}\left(\frac{1^{x}+2^{x}+\cdots+2026^{x}}{2026}\right)^{\frac{1}{\sin x}}$\par
% p079-b19; source PDF page 79
\blockanchor{p079-b19}
\question{3}计算重积分 $\iint_{D} (x^{2} + y^{2}) \, \mathrm{d}x \, \mathrm{d}y$, 其中 $D$ 由椭圆 $9x^{2} + y^{2} = 1$ 所围成.\par
% p079-b20; source PDF page 79
\blockanchor{p079-b20}
\question{4}想不起来了.\par
% p079-b21; source PDF page 79
\blockanchor{p079-b21}
\question{5}计算定积分 $\int_{-1}^{1} \frac{t^{2027} - t^{2025} + t + 1}{\cos^2 t} \, \mathrm{d}t.$\par
% p079-b22; source PDF page 79
\par\Needspace{5\baselineskip}
\blockanchor{p079-b22}
\majorquestion{二}(15 分) 计算曲线积分 $\oint_{L} \frac{x \, \mathrm{d}y - y \, \mathrm{d}x}{16x^{2} + y^{2}}$, 其中 $\Gamma$ 为圆周 $(x - 2)^{2} + y^{2} = 9$, 取逆时针方向.\par
% p079-b23; source PDF page 79
\par\Needspace{5\baselineskip}
\par\Needspace{6\baselineskip}
\blockanchor{p079-b23}
\majorquestion{三}(15 分) 设函数 $u(x, y)$ 由方程组\par
\nopagebreak[4]
% p079-b24; source PDF page 79
\blockanchor{p079-b24}
\displayformula{\left\{ \begin{array}{l} y z + 2 x z + x y = u, \\ x ^ {2} + y ^ {2} + z ^ {2} = 1 \end{array} \right.}
% p079-b25; source PDF page 79
\blockanchor{p079-b25}
所确定, 求  $\frac{\partial u}{\partial x}$  及  $\frac{\partial^{2}u}{\partial y\partial x}$ .\par
% p080-b00; source PDF page 80
\par\Needspace{5\baselineskip}
\par\Needspace{6\baselineskip}
\blockanchor{p080-b00}
\majorquestion{四}(15 分) 设 $k(x, y) \in C([0, 1] \times [0, 1])$, $f(x)$ 在 $[0, 1]$ 上可积, 且\par
\nopagebreak[4]
% p080-b01; source PDF page 80
\blockanchor{p080-b01}
\displayformula{h (x) = \int_ {0} ^ {x} f (y) k (x, y) \mathrm{d} y.}
% p080-b02; source PDF page 80
\blockanchor{p080-b02}
证明： $h(x)$ 在[0,1]上连续.\par
% p080-b03; source PDF page 80
\par\Needspace{5\baselineskip}
\blockanchor{p080-b03}
\majorquestion{五}(15 分) 求由曲线 $y = 2x^{2}$ 与 $y = x$ 所围成的平面图形的面积, 并求\par
% p080-b04; source PDF page 80
\blockanchor{p080-b04}
(1) 该图形绕 x 轴旋转所得的旋转体体积.\par
% p080-b05; source PDF page 80
\blockanchor{p080-b05}
(2) 该图形绕 $y$ 轴旋转所得的旋转体体积.\par
% p080-b06; source PDF page 80
\par\Needspace{5\baselineskip}
\blockanchor{p080-b06}
\majorquestion{六}（15分）求函数 $f(x,y) = x^{2} + y^{2} - 6x + 8y$ 在区域 $D:x^{2} + y^{2}\leq 25$ 上的最大值与最小值.\par
% p080-b07; source PDF page 80
\par\Needspace{5\baselineskip}
\blockanchor{p080-b07}
\majorquestion{七}(20 分) 设函数列 $f_{n}(x) = x^{2025n} - x^{2026n}$.\par
% p080-b08; source PDF page 80
\blockanchor{p080-b08}
(1) 证明函数列 $\{f_n(x)\}$ 在 $\left[0, \frac{2025}{2026}\right]$ 上一致收敛.\par
% p080-b09; source PDF page 80
\blockanchor{p080-b09}
(2) 说明函数列 $\{f_n(x)\}$ 在 $[0,1]$ 上不一致收敛.\par
% p080-b10; source PDF page 80
\par\Needspace{5\baselineskip}
\par\Needspace{6\baselineskip}
\blockanchor{p080-b10}
\majorquestion{八}(15 分) 设 $f(x)$ 在 $[0, 1]$ 上连续, 在 $(0, 1)$ 上可导, 且\par
\nopagebreak[4]
% p080-b11; source PDF page 80
\blockanchor{p080-b11}
\displayformula{f ^ {2} (1) - f ^ {2} (0) = 2026.}
% p080-b12; source PDF page 80
\blockanchor{p080-b12}
证明：存在 $\xi \in (0,1)$，使得 $f'(\xi)f(\xi) = 2026\xi$\par
% p080-b13; source PDF page 80
\par\Needspace{5\baselineskip}
\blockanchor{p080-b13}
\majorquestion{九}(10 分) 设正数列 $\{a_{n}\}$ 满足 $a_{n+1} \leq a_{n} + \frac{1}{2026^{n}}$.\par
% p080-b14; source PDF page 80
\blockanchor{p080-b14}
(1) 证明数列 $\{a_{n}\}$ 有界.\par
% p080-b15; source PDF page 80
\blockanchor{p080-b15}
(2) 证明数列 $\{a_{n}\}$ 收敛.\par
